\subsectionaddtolist{الف}

\subsubsection*{بلوک افزایش نمونه (\lr{Upsampler} L $\uparrow$):}
خروجی برابر $g[n] = s[n/L]$ برای $n$های مضرب $L$ و در غیر این صورت صفر است. تبدیل فوریه را مستقیما محاسبه می‌کنیم؛ تنها جملات مضرب $L$ ناصفرند، پس با تغییر متغیر $n = kL$:
\[
G(e^{j\omega}) = \sum_{n=-\infty}^{\infty} g[n]\,e^{-j\omega n}
              = \sum_{k=-\infty}^{\infty} s[k]\,e^{-j\omega k L}
              = S(e^{j\omega L})
\]
یعنی طیف ورودی به اندازه‌ی عامل $L$ در محور فرکانس فشرده می‌شود و در بازه‌ی $[0,2\pi)$ به تعداد $L$ نسخه (تصویر یا \lr{Image}) از طیف اصلی ظاهر می‌شود.

\subsubsection*{شرط عدم اعوجاج:} 
چون $S(e^{j\omega})$ متناوب با دوره‌ی $2\pi$ است، $S(e^{j\omega L})$ متناوب با دوره‌ی $2\pi/L$ می‌شود و نسخه‌ها هیچ‌گاه روی هم نمی‌افتند. بنابراین افزایش نمونه به‌خودی‌خود همواره برگشت‌پذیر است و اعوجاجی (پدیده‌ی \lr{Aliasing}) ایجاد نمی‌کند و شرط خاصی بر ورودی لازم نیست؛ تنها برای حذف تصویرها و بازسازی سیگنال میان‌یابی‌شده باید یک فیلتر پایین‌گذر ایده‌آل با فرکانس قطع $\pi/L$ پس از این بلوک قرار گیرد.

\subsubsection*{بلوک کاهش نمونه (\lr{Downsampler} $\downarrow M$):}
خروجی برابر $g[n] = s[nM]$ است. تبدیل فوریه‌ی آن از مجموع نسخه‌های جابه‌جا و کشیده‌شده‌ی طیف ورودی به دست می‌آید:
\[
G(e^{j\omega}) = \frac{1}{M}\sum_{k=0}^{M-1} S\!\left(e^{\,j\frac{\omega - 2\pi k}{M}}\right)
\]
طیف ورودی با عامل $M$ کشیده شده و $M$ نسخه‌ی جابه‌جاشده با هم جمع می‌شوند.

\subsubsection*{شرط عدم اعوجاج:}
 برای آنکه این $M$ نسخه روی هم نیفتند (عدم \lr{Aliasing})، باید ورودی باندمحدود باشد:
\[
S(e^{j\omega}) = 0 \qquad \text{برای}\quad \frac{\pi}{M} \leq |\omega| \leq \pi
\]
یعنی پهنای باند سیگنال ورودی باید کمتر از $\pi/M$ باشد؛ در غیر این صورت معمولا پیش از کاهش نمونه یک فیلتر پایین‌گذر ضداعوجاج با فرکانس قطع $\pi/M$ به کار می‌رود.

\subsectionaddtolist{ب}

می‌خواهیم نشان دهیم سیستم \lr{(a)} (ابتدا $\downarrow M$ سپس $H(z)$) با سیستم \lr{(b)} (ابتدا $H(z^M)$ سپس $\downarrow M$) معادل است. رابطه‌ی کاهش نمونه در حوزه‌ی $z$ را با تعریف $W_M = e^{-j2\pi/M}$ به‌صورت زیر می‌نویسیم:
\[
\big[\downarrow M\big]: \quad Y(z) = \frac{1}{M}\sum_{k=0}^{M-1} X\!\left(z^{1/M} W_M^{\,k}\right)
\]

{سیستم \lr{(a)}:} ابتدا کاهش نمونه و سپس فیلتر:
\[
Y_a(z) = H(z)\cdot X_a(z)
       = H(z)\cdot \frac{1}{M}\sum_{k=0}^{M-1} X\!\left(z^{1/M} W_M^{\,k}\right)
\]

{سیستم \lr{(b)}:} ابتدا فیلتر $H(z^M)$ که خروجی‌اش $X_b(z) = H(z^M)X(z)$ است، سپس کاهش نمونه:
\[
Y_b(z) = \frac{1}{M}\sum_{k=0}^{M-1} X_b\!\left(z^{1/M} W_M^{\,k}\right)
       = \frac{1}{M}\sum_{k=0}^{M-1} H\!\left(\big(z^{1/M} W_M^{\,k}\big)^{M}\right) X\!\left(z^{1/M} W_M^{\,k}\right)
\]
اکنون توان $M$ را ساده می‌کنیم. چون $W_M^{\,kM} = e^{-j2\pi k} = 1$، داریم:
\[
H\!\left(z\, W_M^{\,kM}\right) = H(z)
\]
بنابراین $H$ از مجموع بیرون می‌آید:
\[
Y_b(z) = H(z)\cdot \frac{1}{M}\sum_{k=0}^{M-1} X\!\left(z^{1/M} W_M^{\,k}\right) = Y_a(z)
\]
پس دو سیستم معادل‌اند.

\subsectionaddtolist{ج}

می‌خواهیم نشان دهیم سیستم \lr{(a)} (ابتدا $H(z)$ سپس $\uparrow L$) با سیستم \lr{(b)} (ابتدا $\uparrow L$ سپس $H(z^L)$) معادل است. رابطه‌ی افزایش نمونه در حوزه‌ی $z$ از قسمت اول نتیجه می‌شود:
\[
\big[\uparrow L\big]: \quad \text{اگر ورودی } U(z) \text{ باشد، خروجی } U(z^L) \text{ است.}
\]

{سیستم \lr{(a)}:} ابتدا فیلتر، خروجی $X_a(z) = H(z)X(z)$، سپس افزایش نمونه:
\[
Y_a(z) = X_a(z^L) = H(z^L)\,X(z^L)
\]

{سیستم \lr{(b)}:} ابتدا افزایش نمونه، خروجی $X_b(z) = X(z^L)$، سپس عبور از فیلتر $H(z^L)$:
\[
Y_b(z) = H(z^L)\,X_b(z) = H(z^L)\,X(z^L)
\]
چون $Y_a(z) = Y_b(z) = H(z^L)X(z^L)$، دو سیستم معادل‌اند.
